% !TeX root = ../教材题库_全部条目_转写工作稿.tex
% Source: XJ-B1-C06-HW-P256-Y1125
% 第6章 / 6.4 用样本估计总体 / 教材 p.256、257 / PDF 第263、264页
% OCR未校对：公式与文字均可能有误；未完成数学TeX转写；禁止直接发给学生
% 原题及相邻原解截图（核对用）：../题目原图/XJ-B1-C06-HW-P256-Y1125-01.webp, ../题目原图/XJ-B1-C06-HW-P256-Y1125-02.webp
% 本文件只在转写完成后作为题干使用；OCR 条目尚不是可靠题干。
\begin{bankitem}{习题6.4 第11题}{第6章 / 6.4 用样本估计总体 / 教材 p.256、257 / PDF 第263、264页}{XJ-B1-C06-HW-P256-Y1125}{OCR未校对：公式与文字均可能有误；未完成数学TeX转写；禁止直接发给学生}
\begin{ocrdraft}
[教材p.256 / PDF第263页]
\par
11.如图是一个景区甲、乙两段台阶路，图中所标数字表示台阶的高度（单
\par
位：(m).试结合所学的统计方法（如运用平均数、中位数、方差和极差）来评价走
\par
哪段台阶路更舒月廴
\par
15
\par
19
\par
14
\par
10
\par
14
\par
17
\par
16
\par
18
\par
16
\par
15
\par
11
\par
甲
\par
乙
\par
（第11题）
\par
[教材p.257 / PDF第264页]
\end{ocrdraft}
\end{bankitem}
